\section{Case studies}\label{sec:cases}

Three studies follow, each shipping as a runnable example directory and each a
paired experiment: runs differing by exactly one configuration line. Because
every random draw comes from a stream keyed by name rather than call order,
every other class's arrivals, sizes and durations stay byte-identical across
the pair.

\subsection{Frontier scale: running hot without starving
production}\label{sec:case-frontier}

\texttt{examples/04\_frontier/} is 4 clusters $\times$ 32 pods $\times$ 16
scalable units $\times$ 32 nodes $\times$ 8 chips = 524,288 H100 GPUs, run for
two simulated days with a 120\,s scheduler round at seed 42. Node failures are
off deliberately: a 16,384-node gang is eight times the size of Llama~3's 405B
run, which measured one unplanned interruption every 3.1 hours \cite{llama3}, so
it would break about every 23 minutes (\S\ref{sec:failures} gives the
derivation and the disagreement between the two published anchors), and the
restart churn would drown the scheduling effect. This is therefore a
failure-free result. The two runs differ only in whether a best-effort class
exists: the lowest, always-preemptible priority band, in the sense Borg made
standard \cite{borg2015}. A closed loop holds a fixed number of those jobs in
flight; over the two-day horizon that cycles 11,252 of them through the fleet.

\begin{table}[t]
\centering
\begin{tabular}{lrr}
\toprule
 & No backlog & With backlog \\
\midrule
Occupancy (window) & 0.477 & 0.995 \\
Goodput (window) & 0.966 & 0.880 \\
Productive fraction, occupancy $\times$ goodput & 0.461 & 0.876 \\
\midrule
Eval queue wait, p50 & 61.3\,s & 77.8\,s \\
Batch fine-tune queue wait, p50 & 66.9\,s & 116.4\,s \\
Pretrain queue wait, p50 & 51.7\,s & 163.8\,s \\
Eval queue wait, p99 & 236\,s & 1{,}292\,s \\
Preemptions per minute & 0.068 & 0.565 \\
Jobs finished & 23{,}538 & 34{,}422 \\
\bottomrule
\end{tabular}
\caption{A 524,288-chip fleet with and without a preemptible backlog, seed 42.
The first three rows are measured over the steady-state window (the middle
80\,\% of the run); the rest are full-run figures. The two scopes are not
interchangeable, which is why the example labels them.}
\label{tab:frontier}
\end{table}

Table~\ref{tab:frontier} is the trade. The backlog roughly doubles
steady-state occupancy and finishes 10,884 more jobs. Production medians move
by one or two scheduler rounds --- the largest is the pretrain median, 51.7\,s
to 163.8\,s --- which is the sense in which production waits stay roughly
flat; the tails do not, the eval p99 going from 236\,s to 1,292\,s, roughly 4
minutes to 22. The cost is goodput, 8.6 points of it: the backlog runs with
checkpointing off, which is exactly what makes it instantly reclaimable, so
every eviction discards a whole \emph{stint} --- one contiguous period of
holding an allocation, and with no checkpoint behind it, everything that job
computed since it was last given chips. The product of the two
still nearly doubles. Best-effort jobs are excluded from every wait
distribution by contract, waiting time being meaningless for a closed loop at
saturation.

Then one large job: a 131,072-chip pretraining gang of 32 pod-sized segments,
modelled on the multi-pod shape Llama-3 405B training used \cite{llama3},
submitted at hour 20.57 of both
runs. On the idle fleet it waits 113\,s; above 99\,\% occupancy, 233\,s ---
exactly one 120\,s round more, so the reclaim costs nothing and the wake
boundary costs everything. One wake carries 345 preemptions, allocated chips
dip from 523,946 to 406,406 and recover to 522,998 a round later, and pending
jobs spike from 116 to 441. The gang spans 32 pods in both runs.

\subsection{Occupancy and goodput can move in opposite
directions}\label{sec:case-tradeoff}

\texttt{examples/05\_topology\_tradeoff/} is 8,192 H100s (8 pods $\times$ 16
racks $\times$ 8 nodes $\times$ 8 chips) over seven days at seed 42, failures
off, near 86\,\% utilization so pod-shaped holes are scarce. Pretraining jobs
draw 256, 512 or 1,024 chips and ask for one pod, but softly: they wait ten
minutes for a pod-shaped hole, then take whatever fits. Run~A charges a job
spanning more than one pod a speed multiplier of 0.7, so it needs $1/0.7
\approx 1.43$ times the wall time for the same work; run~B sets that
multiplier to 1.0.

\begin{table}[t]
\centering
\begin{tabular}{lrr}
\toprule
 & A: crossing costs 0.7$\times$ & B: crossings free \\
\midrule
Occupancy (window) & 0.862 & 0.826 \\
Goodput (window) & 0.861 & 0.983 \\
\midrule
Cross-pod placements & 35 of 112 & 17 of 112 \\
Pretrain queue wait, p50 & 484.5\,s & 47.9\,s \\
Pretrain job completion time, p50 & 51{,}002\,s & 43{,}136\,s \\
Pretrains completed & 99 & 103 \\
Eval queue wait, p50 & 31.9\,s & 31.4\,s \\
\bottomrule
\end{tabular}
\caption{Turning the cross-pod speed penalty on raises occupancy and lowers
goodput at the same time (seed 42, 7-day horizon).}
\label{tab:penalty}
\end{table}

Switching the penalty \emph{on} raises occupancy, 0.826 to 0.862, and
\emph{lowers} goodput, 0.983 to 0.861. Both are correct: occupancy is the
fraction of healthy chip-time allocated to some job, goodput the fraction of
that allocated time which produced surviving work. A slowed job holds its
chips about 1.43 times longer for the same work, inflating the first and
deflating the second. The busier-looking run completed four fewer pretrains
and finished them later.

Table~\ref{tab:penalty} also shows the feedback loop that makes this
self-reinforcing, and it reads backwards at first: charging for pod crossings
\emph{doubles} them, from 17 of 112 placements to 35. The chain is that a
penalised job holds its chips 1.43 times longer, so pod-shaped holes are
scarcer when the next pretraining job arrives, so more jobs exhaust their
ten-minute soft wait and take whatever fits. Pretrain queue wait moves
accordingly, 47.9\,s to 484.5\,s.

Two details keep this honest. Run~B's goodput is 0.983 rather than 1.0 because
of checkpoint-save amortization --- a one-hour interval with a 60\,s save
costs about 1.6\,\% --- so the penalty accounts for 12 of those points, not
14. And the effect is confined to large jobs: eval waits are 31.9\,s and
31.4\,s, while the chance of crossing pods rises from 18\,\% at 256 chips to
61\,\% at 1,024. This is the strongest argument available against an
occupancy-only dashboard: occupancy rose 3.6 points and the fleet got worse.

\subsection{Placement: where the small jobs land decides how long the big ones
wait}\label{sec:case-placement}

\texttt{examples/07\_placement\_study/} is deliberately tiny --- 4 racks
$\times$ 8 nodes $\times$ 8 chips = 256 H100s in 32 nodes --- and deliberately
two levels deep, since \texttt{best\_fit} and \texttt{consolidate} are
bit-identical on a single-level fleet. Fourteen days, seed 42, failures off.
The workload is mice (1-, 2- and 4-chip jobs lasting hours) plus gangs of
16, 32 and 64 chips --- 2, 4 and 8 \emph{whole nodes} --- at an offered load of
78\,\% by construction, 81\,\% realised at this seed, because Poisson arrivals
and a lognormal duration tail land above the mean here. A best-effort FIFO
scheduler holds ordering fixed, so only the
placement policy varies. It reproduces the mechanism the Helios replay
identified: a 1-chip job landing on a completely free node converts eight
whole-node chips into a remainder no gang can claim, and the gangs, not the
mice, pay for it.

\begin{table}[t]
\centering
\begin{tabular}{lrrrr}
\toprule
 & \texttt{first\_fit} & \texttt{best\_fit} & \texttt{consolidate} & \texttt{spread} \\
\midrule
Stranded whole nodes, mean (of 32) & 6.55 & 5.31 & 5.37 & 25.31 \\
64-chip gang queue wait, mean (s) & 19{,}056.8 & 16{,}195.3 & 16{,}314.7 & 739{,}842.1 \\
64-chip gangs ever started (of 109) & 107 & 107 & 107 & 29 \\
Mice queue wait, p50 (s) & 31.0 & 31.4 & 31.3 & 30.1 \\
Occupancy (full run) & 0.8141 & 0.8141 & 0.8141 & 0.5829 \\
Jobs completed & 3{,}436 & 3{,}436 & 3{,}436 & 3{,}309 \\
\bottomrule
\end{tabular}
\caption{Four placement policies on identical hardware, identical arrivals and
identical ordering (seed 42, 14-day horizon). A ``stranded'' node is a healthy
node that is partly occupied, so its free chips cannot satisfy any whole-node
request. The whole study runs in about 10\,s on a laptop: about 1\,s each for
the three packing arms and about 7\,s for \texttt{spread}, whose queue is
roughly 30$\times$ deeper.}
\label{tab:placement}
\end{table}

Tightest-fit packing (\texttt{best\_fit}; \texttt{consolidate} is within noise
of it here) cuts time-average stranded nodes from 6.55 to 5.31 of 32
and the mean queue wait of an 8-node gang from 19,056.8\,s to 16,195.3\,s
(Table~\ref{tab:placement}). The mice pay for it: their median wait is flat,
31.0 against 31.4\,s, but their mean rises from 80.7\,s to 95.3\,s --- about
15 seconds of mouse wait buying about 2,900 seconds of gang wait.
\texttt{spread} prices the mechanism from the other side: it strands 25.31 of
32 nodes, never starts 80 of the 109 8-node gangs and delivers 28.4\,\% fewer
chip-hours on identical hardware --- and it averages only the gangs it did
start.

The negative result matters more. Full-run occupancy, goodput and allocated
chip-hours are \emph{bit-identical} across \texttt{first\_fit},
\texttt{best\_fit} and \texttt{consolidate} --- 0.8140502977503281,
0.9831464030337717 and 70,021.35041129222 chip-hours --- because the same
3,436 jobs complete for the same durations in all three, and windowed
occupancy moves 0.2 percentage points. The entire placement effect is
\emph{when} work ran, never how much: a placement study reporting only
occupancy has measured almost nothing.

Three caveats travel with it. At seed 1234 the mechanism metric still improves
by about 20\,\% while the outcome regresses: mean gang wait 8,205\,s to
9,407\,s, 15\,\% worse. A hard \texttt{within: rack} constraint inverts the
ranking outright --- 64-chip mean wait 39,677\,s under \texttt{first\_fit}
against 121,486\,s under \texttt{best\_fit} --- because node-level
tightest-fit spreads gangs across racks, so no rack ever drains. And this
scheduler is best-effort, so a stuck gang does not hold up the mice behind it:
a floor, not a best case. ``\texttt{best\_fit} is a better scheduler'' is not
a claim this study supports.

\section{Implementation, testing and reproducibility}\label{sec:impl}

\paragraph{Language and size.} fleetsim is Python 3.11 or newer under
Apache-2.0, depending at runtime on numpy, pandas, pyarrow, pyyaml and
optionally matplotlib. The event loop is hand-written rather than built on a
simulation framework --- full control over tie-breaking and determinism is
worth about 200 lines --- and the local web application is pure
standard-library \texttt{http.server}, so running it adds no dependency. The
simulator is 22,563 lines of Python across 47 files; the tests are 18,896
lines across 36 files plus 2,585 lines of validation rungs across 12, and
pytest collects 1,024 of them. The four placement policies behind
Section~\ref{sec:case-placement} are 340 lines together; the FIFO scheduler is
48.

\paragraph{Determinism, and its per-platform limit.} Simulated time is a
64-bit integer count of microseconds, so ordering and tie-breaking are exact.
Random draws come from numpy seed sequences whose spawn key is the first 16
bytes of the SHA-256 of the stream's \emph{name}, never Python's salted
\texttt{hash()}, so a
given stream yields the same numbers in every process on every machine. The
contract is that an identical scenario and seed produce byte-identical Parquet
and JSON, on a given platform, and a CI test enforces it by running the whole
pipeline twice with failures, maintenance and preemption switched on.

That qualifier is load-bearing. The lognormal and exponential samplers call
the platform's C math library, whose \texttt{exp} and \texttt{log} differ in
the last bit between glibc and Apple's implementation. An early release (v0.1;
the project is at v0.8) shipped a compatibility test pinning a SHA-256 of the
arrival stream captured on macOS, and it failed the first time CI ran on Ubuntu. The replacement compares golden
\emph{data} --- the first 1,000 arrivals at nine significant digits, plus
aggregate counts --- with tolerances: 2 microseconds on times, $10^{-6}$
relative on durations, and at most 5 differing rows per 1,000 for draws
sitting on a discrete knife edge. Cross-platform reproduction means the same
distributions, not the same bytes.

\paragraph{Continuous integration, and what a clean machine caught.} One
GitHub Actions workflow runs on every push and pull request: Ubuntu, a Python
3.11 and 3.12 matrix, an editable install of the package and the out-of-tree
plugin example, \texttt{pytest tests validation -q}, then a smoke test running
\texttt{fleetsim validate} and \texttt{fleetsim run} on the minimal example.
It runs with no network in about 80\,s with exactly four skips --- two Helios
full-trace rungs, one Philly, one frontier-budget gate --- and anything else
skipping is a signal.

Two defects surfaced only because CI is a clean room. The first was an
environment leak: a test asserting that every shipped example validates passed
locally only because the developer's virtual environment still had the
example-03 plugin scheduler installed from an earlier session, and that
example names an out-of-tree scheduler which enters the registry only when its
package is installed. The fix skips an example whose scheduler is unregistered
--- validation \emph{should} reject unknown names --- while asserting a floor
on how many examples were actually checked, so the skip cannot hollow the test
out; CI now installs the plugin. The second was a concurrency deadlock a
laptop could not reach: the 3.12 job hung at 91\,\% and was killed by the job
timeout while 3.11 passed the same commit, the developer machine having many
cores and the hosted runners two. The cause was a helper process that starts
lazily, combined with a setup step that re-ran on every worker-pool creation and
mutated interpreter-global state; on a two-core runner the two raced. Building
that context once per process fixed the hang, and a faulthandler timeout was
added so that the next hang of any kind would fail loudly with every thread's
stack. That dump then named a deeper bug: CPython invokes a future's
done-callback while holding the executor's shutdown lock, so a callback
calling \texttt{shutdown()} deadlocks itself. A broken pool is now disposed on
a separate daemon thread.

\paragraph{Adversarial review.} Several releases were reviewed adversarially ---
a pass whose goal was to find defects rather than approve --- and each such pass
leaves a regression file that lists its findings:
\texttt{test\_\allowbreak review\_\allowbreak fixes.py} in \texttt{tests/}, plus
its \texttt{v02}, \texttt{v04}, \texttt{v05} and \texttt{v05\_hardening}
successors. Five files across eight releases, so this is a habit rather than a
gate. The same
discipline governs the published figures: each lives once, as data, in a
single module from which the validation tests import their assertion constants
and which the web application serves verbatim, and a test checks by object
identity that nobody has retyped a number into a test.

\paragraph{Reproducing the validation.} The offline checks are one command,
\texttt{pytest tests validation -q}, and \texttt{fleetsim validation run}
prints the vendored-slice checks in about 5\,s. The trace results are opt-in
and environment-gated because they download real data:
\texttt{FLEETSIM\_HELIOS\_FULL=1} \texttt{pytest -m trace\_full}
\texttt{validation/test\_helios\_ratio.py} fetches the 36\,MB Helios archive
\cite{helios2021} and replays four clusters under two schedulers in 266\,s.
The Philly rung \cite{philly2019} is gated the same way but needs \texttt{git
lfs pull} first, since a plain fetch of that 1\,GB artifact yields a 135-byte
pointer; it was written to specification and has not been run against the real
trace in this build. Each rung self-skips unless its variable is set and never
downloads implicitly. \texttt{fleetsim validation cite} prints each trace's
required attribution.

\paragraph{Tooling.} Six subcommands cover the loop: \texttt{fleetsim
validate}, \texttt{run}, \texttt{compare}, \texttt{plot}, \texttt{viz} and
\texttt{serve}. Two properties of them carry an argument rather than a feature
list. \texttt{viz} renders a run directory into one self-contained HTML replay
that makes zero external requests, and a test greps both the template and the
output for any external URL or fetch call, so the claim is enforced rather than
intended. And parameter sweeps run one simulation per grid cell in worker
processes: the same sweep run twice is byte-identical cell for cell, so
parallelism does not weaken determinism, and a test enforces that too.
